\section*{People Power \& Conflict Evolution}

There is little empirical evidence that protests should or should not ``work" in contexts of high violence at the sub-set level. At present, several possible logics linking the protests and violence between armed groups exist. 

First, it is possible that there is no relationship. Protest might have other effects and thus persist for reasons outside of the stated goals driving the protest organizers. Protests might build community between survivors of violence and foster a sense of purpose and belonging in the midst of crisis. Protests against insecurity and issues of violence might be targeted and both types of armed actors (i.e. the state and the non-state challenger) these actors might respond through other means outside of their violent tactics. Non-state groups might make public appeals to civilians, saying that they are there for civilians' protections. Similarly, the government can respond with media campaigns, speeches, or generally try to divert political attention from violent events. A second mechanism driving a null relationship between protests and violent events is the high level of risk for protestors. Protestors might be targeted because of their activism, possible resulting in fewer protests over time but with no clear consequences for violence between armed groups.

Second, protests might achieve their stated aims and decrease violent contestation in their region. If protests successfully demand that the government invest greater resources in the prevention of violence, then we would expect fewer incidents of violence following protests. If protests also demand that non-state actors stop warfare with the state, this becomes a public disapproval of any attempt by the non-state group to coerce or control the civilian population. With enough discontent, protests disrupt the ability of the non-state actor to operate in their region.

As violence between armed actors (state and non-state) increases, non-state actors often victimize civilians to highlight the government's inability to protect the population. Even if protestors demand a non-violent solution to warfare, a final plausible logic connecting protests and violence is that protest raises the cost of government inaction, so that government armed actors are more likely visibly fight hard in the next round of contestation. 

